\documentclass{article}

\title{Apa yang membawaku ke RSC Aksantara ITB}
\author{Ivan Ahmadi H., 16525136 / 13225111}

\begin{document}
\maketitle

Sebagai mahasiswa STEI-Rekayasa ITB, aku merasa bahwa terdapat banyak wadah di ITB yang dapat memumpuni minat-minatku, terutama di bidang sains dan teknologi.
Dari dulu, aku senang sekali meng-eksplorasi tentang robot, dan robotika di dalamnya. Saat SD, aku sempat mengikuti les robotik yang membuka mataku terhadap dunia baru, serta menantang cara berpikir dan logikaku. 
Les robotik ini yang akhirnya membawaku ke ITB dan meneruskan pendidikan tinggi di program studi Teknik Elektro. 
Ketertarikanku terhadap aksantara berawal dari rekomendasi kakakku, beliau juga merupakan alumni teknik elektro ITB yang memperkenalkanku kepada aksantara. 

Aksantara merupakan salah satu UKM yang sejak awal masuk ITB aku tertarik untuk ikut serta meskipun belum mengetahui banyak mengenai UKM ini, hingga saat OHKM aku mencoba untuk menggali informasi mengenai aksantara, dimana pada saat itu aku mendapatkan informasi bahwa aksantara baru akan memulai Open Recruitment (OPREC) di semester genap.
Aku pun menjalankan semester pertamaku di ITB dengan antisipasi menunggu OPREC aksantara dibuka, sekaligus mendalami kegiatan akademis dan organisasional.
Ketika aku mendapatkan informasi bahwa aksantara sudah memulai OPREC, aku langsung mendaftar tanpa berpikir panjang, dan aku pun menunggu informasi mengenai OPREC selanjutnya.
Pada tanggal 16 Januari 2026, informasi pertama mengenai OPREC diumumkan, yaitu Online Onboarding lewat MS Teams.
Aku pun akhrnya mengikuti Online Onboarding, lalu dikenalkan dengan beberapa departemen di aksantara, sekaligus teknis untuk bagaimana cara melanjutkan OPREC dan daftar ke departemen tersebut.
Pada saat Online Onboarding, 2 departemen langsung menarik perhatianku, yaitu departemen Electrical System dan RSC.

Setelah Online Onboarding, aku pun melihat kriteria dan syarat untuk lanjut mengikuti OPREC, salah satu dari syarat tersebut merupakan THT/Take Home Test yang berbeda untuk setiap depertemennya.
Aku langsung melihat THT dari Electrical System dan RSC, dan akhirnya memutuskan untuk mengambil THT dari RSC, karena menurutku aku bisa eksplorasi dalam di RSC, dan skill pemrograman yang dibutuhkan di RSC dapat menjadi bekal pembelajaran yang dapat membantuku di pendidikanku seterusnya.
Dalam pengerjaan THT RSC, satu hal sangat menantang bagiku, yaitu penggunaan GitHub sebagai media utama dalam pengerjaan dan pengumpulan THT.
Akhirnya aku pun mempelajari GitHub secara mendalam, dan pengeraan THT ke depannya berjalan sesuai harapan. Alhamdulillah pada saat pengumuman pendidikan, aku lolos seleksi dengan nilai yang cukup memuaskan.
Setelah lolos seleksi, aku pun mengikuti pendidikan RSC yang sedang berjalan selama kurang lebih satu minggu ini, ilmu yang didapatkan sejauh ini cukup menarik, bermanfaat, dan menantang.
Tugas-tugas yang tidak mudah sama sekali, membuatku semakin termotivasi untuk belajar dan meneruskan pendidikan ini sampai selesai.

Aku sangat bersyukur bisa mendapatkan kesempatan untuk mengikuti pendidikan RSC Aksantara ITB ini, dan apabila diberi kesempatan lanjut untuk menjadi anggota resmi dari aksantara ITB, tentu saja aku akan aktif menjalankan peranku sebagai RSC, baik dimanapun tim aku ditempatkan.
Aku berharap dengan menjadi anggota RSC Aksantara ITB, aku dapat menjadi anggota yang bermanfaat bagi aksantara, dan terus belajar mendalami ilmu ini, serta mengaplikasikannya di dunia nyata.




\end{document}